\documentclass[12pt, oneside]{book}

\usepackage{graphicx}
\usepackage{amsmath}
\usepackage{amsfonts}
\usepackage{amssymb}
\usepackage{hyperref}

\usepackage[utf8]{inputenc}
\usepackage[T2A,T1]{fontenc}
\usepackage{ragged2e}

\usepackage{setspace}
\onehalfspacing
\setlength{\parindent}{1.25cm}

\usepackage[english,russian]{babel}
\usepackage[a4paper, left=3.5cm, top=2cm, right=2.5cm, bottom=2cm]{geometry}

\makeatletter
\renewcommand{\chapter}{%
        \thispagestyle{plain}
        \global\@topnum\z@
        \@afterindentfalse
\secdef\@chapter\@schapter}

\renewcommand{\@makechapterhead}[1]{%
        {\parindent \z@ \raggedright
                \normalfont
                \LARGE \bfseries \thechapter\quad #1\par\nobreak
                \vskip 40\p@
}}
\makeatother

\begin{document}
\pagestyle{plain}
\thispagestyle{empty}
\begin{center}
        \textbf{\large САНКТ-ПЕТЕРБУРГСКИЙ ГОСУДАРСТВЕННЫЙ УНИВЕРСИТЕТ} \\ [4cm]

        {\large Отчёт по курсу <<Современные проблемы прикладной математики>>} \\ [2cm]

        \textbf{\Large Математическое моделирование и реализация
                планировщика процессов для гетерогенных
                процессорных архитектур} \\ [2cm]

        {\large Шаго Павел Евгеньевич} \\ [0.3cm]
        {22.Б07-ПУ, 01.03.02 Прикладная математика и информатика} \\ [3cm]

        {Проверил} \\ [0.3cm]
        {кандидат ф.-м.н., доцент} \\ [0.3cm]
        {Кривошеин Александр Владимирович} \\ [4cm]

        {Санкт-Петербург} \\
        {20 апреля 2026 г.}
\end{center}


\newpage
\chapter*{Постановка задачи}

\quad Современные процессорные архитектуры всё чаще объединяют в
одном кристалле ядра с принципиально разными характеристиками
производительности и энергопотребления --- такие платформы принято
называть гетерогенными (Intel Hybrid CPU, ARM big.LITTLE).

\quad Стандартный алгоритм планирования Earliest Eligible Virtual
Deadline First (EEVDF) ядра Linux не учитывает эту неоднородность
и выбирает ядра CPU для выполнения задач вне зависимости
от их вычислительной мощности, что может приводить к неоптимальному
использованию ресурсов процессора.

\quad Фреймворк sched\_ext --- подсистема ядра Linux, позволяющая
реализовывать алгоритмы планирования в виде eBPF программы.
Такие программы проходят верификацию и JIT-компиляцию непосредственно
при загрузке и исполняются в изолированной виртуальной машине внутри
ядра Linux; при этом гарантируется безопасность без ущерба производительности.
Ключевое практическое свойство фреймворка состоит в том, что новый
планировщик можно загрузить и выгрузить в работающей системе без
перекомпиляции и перезагрузки ядра, что существенно сокращает
итерационный цикл разработки и экспериментальной проверки
улучшенных подходов.

\quad Объектом исследования являются алгоритмы планирования процессов
в операционных системах на вычислительных устройствах с гетерогенной
процессорной архитектурой.

\quad Целью выпускной работы является разработка и экспериментальная оценка
планировщика, обеспечивающего более высокую эффективность вычислений
на гетерогенном оборудовании по сравнению с базовым
алгоритмом планирования EEVDF ядра Linux.

\newpage
\chapter*{Обзор литературы}

\quad С появлением фреймворка sched\_ext исследовательская активность
в области процессных планировщиков заметно возросла. Большинство
современных работ используют sched\_ext как экспериментальную платформу
для быстрой проверки новых алгоритмов без модификации ядра.
Диапазон применения охватывает оптимизации существующих алгоритмов,
подходы машинного обучения и нестандартные постановки задач.

\quad Tang~et~al.~\cite{srand} предложили sched\_ext планировщик,
использующий облегчённую стратегию SRAND с глобальной очередью в eBPF-map,
добившись сокращения времени переключения контекста на 11,83\%
и роста производительности в стресс-тестах на 7,02\% по сравнению
с EEVDF. Xu~et~al.~\cite{lace} применили sched\_ext в совершенно
иной области --- fuzzing-тестировании параллельного кода ядра:
специализированный CCT-планировщик LACE даёт на 38\% больше покрытия
ветвлений и ускорение исследования в 11,4 раза.
В машинном обучении Mao~et~al.~\cite{aegis} реализовали
планировщик Aegis с управлением на основе DeepQ-Network,
обеспечивающий нулевую потерю событий аудита там, где EEVDF теряет
до 98\% записей.
Wang~et~al.~\cite{mos} показали, что агент на XGBoost,
динамически выбирающий политику из портфеля планировщиков,
превосходит EEVDF в 86,4\% пользовательских сценариев.
Galin и Hedblom~\cite{lavd} оценили планировщик LAVD в
telco-окружении Kubernetes и обнаружили, что минималистичный
scx\_simple при типичных нагрузках превосходит его, несмотря на
значительно меньшую сложность --- результат, подчёркивающий ценность
sched\_ext как платформы быстрого сравнения алгоритмов.

\quad На этом фоне задача планирования на гетерогенных процессорах
остаётся значительно менее изученной.
Наиболее близкой к данному направлению является работа
Van~Craeynest~et~al.~\cite{pie}, в которой показано, что
интенсивность обращений к памяти --- недостаточный критерий
распределения задач по типам ядер. Производительность определяется
совместным влиянием ILP и MLP, что требует оценки на уровне
CPI-компонент. Механизм PIE оценивает производительность на
альтернативном типе ядра без фактической миграции и даёт прирост
пропускной способности на 5,5--8,7\% над конкурирующими подходами.
Однако работа относится к 2012 году, выполнена вне контекста
sched\_ext и рассматривает аппаратный механизм оценки, а не
программную модель виртуального времени.
Именно этот пробел --- учёт вычислительной мощности ядра в рамках
программного планирования на современных гетерогенных платформах и составляет
предмет выпускной работы.

\chapter*{Промежуточные результаты}

\quad Первым этапом реализован планировщик EEVDF~\cite{stoica} на базе
        sched\_ext. Каждой задаче при постановке в очередь назначаются виртуальное
допустимое время $VE_i = \max(v_i,\, V_{\mathrm{now}} - \Delta)$ и
виртуальный дедлайн $VD_i = VE_i + \Delta S / w_i$, где $\Delta$~---
квант sched\_ext, $w_i$~--- вес, $S$~--- масштаб.
Диспетчеризуется задача с наименьшим $VD_i$.

\quad Сравнение с эталонным EEVDF ядра Linux проводилось на трёх
бенчмарках (hackbench, sysbench, sched-latency через eBPF tracepoints).
По пропускной способности реализация превосходит ядерный планировщик из-за
меньшего количества переключений контекста: hackbench показал одинаковое время
0.20~с (+0\%), sysbench вырос с 4157 до 7357~событий/с (+77\%),
что объясняется глобальной очередью sched\_ext, сокращающей число
миграций задач относительно per-CPU очередей ядра.
Повышенный 99-й процентиль задержки (250~мкс против 25~мкс) является
ожидаемым следствием eBPF-прослойки и не значительно сказывается на
быстродействии системы (по hackbench).

\quad Вторым этапом реализован прототип планировщика scx\_auction,
основанного на механизме динамического аукциона Викри (VCG).
В его основе лежит идея о том, что каждый квант процессорного времени
является невозобновляемым ресурсом, а задачи выступают агентами,
конкурирующими за этот ресурс в соответствии со своим приоритетом.
Планировщик различает P-core и E-core как два класса ресурсов с разной
стоимостью кванта ($c_P > c_E$) и разной скоростью исполнения.
Для выбора класса ядра используется коэффициент бюджета задачи
$\rho_i = B_i / B_i^{\max}$. Задача с высоким $\rho_i$ (мало потребляет,
накопила бюджет) расценивается как интерактивная и направляется на P-core,
задача с низким $\rho_i$ (CPU-bound) направляется на E-core.
Бюджет пополняется пропорционально времени простоя и весу задачи,
что создаёт устойчивое равновесие между классами нагрузки.
Задачи с исчерпанным бюджетом получают отложенную диспетчеризацию,
что приближает механизм к инцентивной совместимости теоретического VCG
в условиях бюджетного ограничения.

\newpage
\renewcommand{\bibname}{\Large Список использованных источников}
\begin{thebibliography}{9}
        \vspace*{0.3cm}

	\bibitem{srand}
        Tang~Q., Gao~X., Li~G., Lin~J.
        Optimizing Linux Scheduling Based on Global Runqueue with SCX.
        IEEE International Conference SMC, 2024.

        \bibitem{lace}
        Xu~J., Wolff~D., Yi~H.~X., Li~J., Roychoudhury~A.
        Concurrency Testing in the Linux Kernel via eBPF.
        arXiv, 2025.

        \bibitem{aegis}
        Mao~J., Ujcich~B.~E., Ma~S.
        Rethinking Provenance Completeness with a Learning-Based
        Linux Scheduler.
        arXiv, 2025.

        \bibitem{mos}
        Wang~X., Jia~S., Huang~Z., Cao~J., Song~M.
        Mixture-of-Schedulers: An Adaptive Scheduling Agent as a
        Learned Router for Expert Policies.
        arXiv, 2025.

        \bibitem{lavd}
        Galin~M., Hedblom~A.
        Linux Kernel Scheduler Evaluation for Performance-Critical
        Telecom Workloads.
        Linköping University, 2025.

        \bibitem{pie}
        Van~Craeynest~K., Jaleel~A., Eeckhout~L., Narvaez~P., Emer~J.~S.
        Scheduling Heterogeneous Multi-Cores through Performance Impact
        Estimation (PIE).
        ISCA, ACM, 2012, pp.~213--224.

        \bibitem{stoica}
        Stoica~I., Abdel-Wahab~H.
        Earliest Eligible Virtual Deadline First: A Flexible and Accurate
        Mechanism for Proportional Share Resource Allocation.
        Old Dominion University, 1996.
\end{thebibliography}

\end{document}
